\section{The g-Factor}

A spinning charge in a magnetic field turns at a definite rate. How strongly it turns, for a given amount of spin, is set by one number, the g-factor. For a Dirac fermion the theory says this number is exactly $2$. For the real electron it is measured to be very close to $2$, off only by the small radiative corrections of quantum electrodynamics. The bare value $2$ is one of the cleanest predictions in physics, and any substrate that claims to carry matter must reproduce it without being told to.

This section reads $g = 2$ off the mesh itself. The number is not assigned. It surfaces from the Clifford algebra that the $24$ sites of a dock already carry, through the squared Dirac operator of Section~18. The result fills an old hole, where an earlier test had set the magnetic response by hand. It also draws a sharp line, fixing the \emph{form} of the magnetism while leaving the \emph{strength} of the coupling honestly open.

\begin{principle}[The g-factor is measured, not set]
The spinor's magnetic moment is read off the substrate's own Dirac structure as $g = 2$. It is nowhere put in by hand. This is a genuine measurement at depth L2, with a discriminating control that lacks the moment.
\end{principle}

\subsection{What the g-factor is}

A particle with spin and charge behaves like a tiny magnet. Its magnetic moment $\mu$ is proportional to its spin $S$,
\[
\mu = g \, \frac{q}{2m} \, S,
\]
and the proportionality constant $g$ is the g-factor. A classical spinning ball of charge, with mass and charge spread the same way, gives $g = 1$. The Dirac equation, by contrast, forces $g = 2$. The factor of two is not arbitrary. It is what the Dirac structure produces through the particular way spin couples to the field, and it is the bare value the electron sits next to.

So the question for the mesh is sharp. The $24$ sites of a dock carry spinors, as Section~17 established, with the $D_4$ directions splitting under triality into $8v + 8s + 8c$. If those spinors are genuine Dirac spinors, then placed in a magnetic field they should turn at the Dirac rate, $g = 2$, with no further input. That is what we now check.

\subsection{The old hole}

An earlier magnetism test took a shortcut. It set the gyromagnetic ratio equal to $2$ \emph{by hand}, then integrated the Larmor precession equation and watched a spin precess at the rate that $g = 2$ already implied.

This is circular. If $g = 2$ is assumed at the start, and then a spin is seen precessing at the rate $g = 2$ predicts, nothing about the substrate has been shown. The output is the input wearing a different hat.

\begin{principle}[Assuming g and then measuring g is not a measurement]
A test that sets the gyromagnetic ratio to $2$ and then integrates the precession it implies proves only its own assumption. It is an echo, not a reading. This was logged as an L0 hole, a surface-level demonstration.
\end{principle}

The new result removes the shortcut. Nothing about the magnetic response is assumed. The operator is built, its spectrum is computed, and $g$ is read out of the spectrum.

\subsection{The measurement}

The construction follows the Dirac structure of Section~18 directly. Build the two-component Dirac operator for a tone-knot in a \emph{uniform magnetic field}. The kinetic momenta no longer commute, since a magnetic field twists the two spatial directions into each other. They close instead on the magnetic algebra, which is exactly the algebra of a harmonic oscillator ladder. We carry that algebra exactly, through the oscillator raising and lowering operators, so the field enters the operator with no approximation. Then the operator is diagonalized and its spectrum is read.

What comes out is the \emph{relativistic Landau spectrum}, the ladder of energy levels a Dirac particle takes in a magnetic field. For a plain spinless object the levels are the ordinary Landau rungs, evenly spaced. For the Dirac spinor one feature is different and decisive.

\begin{result}[The zero mode]
The relativistic Landau spectrum of the substrate's Dirac spinor contains a \emph{zero mode}, a level sitting at $E^2 = m^2$, that a spinless scalar in the very same field does not have. This extra level is the spin gap.
\end{result}

The zero mode is the signature of spin. Its presence, and its size, read back directly to the magnetic moment. Working the spin gap of this spectrum through to the g-factor gives $g = 2$.

And $g$ is nowhere put in. It arises from the Clifford algebra that the $24$ sites carry, the same algebra whose square gives the relativistic hinge $H^2 = p^2 + m^2$ of Section~18. The spin term that produces the zero mode is exactly the piece of $H^2$ that a scalar lacks. So the moment is not a separate ingredient. It is what the squared operator already contains.

\begin{principle}[The moment rides on $H^2$]
The spin gap is structural. It is forced by the sites' own Clifford algebra, surfacing through the squared operator $H^2 = p^2 + m^2$. The same hinge that gives the spinor its relativistic dispersion gives it its magnetic moment.
\end{principle}

\subsection{Robustness}

A single number from a single setup could be a coincidence of that setup. What separates a measurement from a coincidence is that the number survives when the conditions change. It does.

\begin{table}[ht]
\centering
\begin{tabular}{@{}lll@{}}
\toprule
varied & values tested & $g$ read back \\
\midrule
field strength & two strengths & $2$ \\
mass & two masses & $2$ \\
\bottomrule
\end{tabular}
\caption{The measured g-factor across varied field strengths and masses. The zero mode sits at $E^2 = m^2$ in every case, and $g = 2$ is stable.}
\label{tab:g-robust}
\end{table}

The same number comes back every time. The zero mode sits at $E^2 = m^2$ regardless of field strength or mass, so $g = 2$ is stable across the conditions tested.

\begin{principle}[Two fields, two masses, one answer]
The g-factor is read as $2$ at two field strengths and two masses. The moment is a property of the sites, not of the setup. A number that holds across the conditions is a measurement, not an artifact.
\end{principle}

\subsection{The control that lacks the moment}

The discriminating control is a \emph{spinless scalar} placed in the very same magnetic field. It feels the field, so it has a Landau spectrum, but it carries no spin, so its spectrum should differ in exactly the spin part.

\begin{table}[ht]
\centering
\begin{tabular}{@{}llll@{}}
\toprule
object & spin term & zero mode & magnetic moment \\
\midrule
Dirac spinor & yes & yes, at $E^2 = m^2$ & $g = 2$ \\
spinless scalar & no & no & none \\
\bottomrule
\end{tabular}
\caption{The spinor against its discriminating control. The scalar's Landau spectrum has no zero mode and no moment. The difference between the two spectra is the spin gap, and the spin gap is $g = 2$.}
\label{tab:g-control}
\end{table}

The scalar's Landau spectrum has no zero mode, no spin gap, and no magnetic moment. It is the test that could have failed and did not. If the zero mode were an artifact of how the field is wired into the operator, the scalar would carry it too. It does not. The gap belongs to the spinor alone.

\begin{principle}[The scalar is the null]
The spinless scalar feels the field but carries no spin, so it shows no gap. The spinor's gap is the difference between the two spectra, and that difference is $g = 2$. The control isolates the moment as a property of spin, not of the field setup.
\end{principle}

\subsection{Why it counts as L2}

The depth rubric is plain. L0 is a surface or circular demonstration. L2 is a genuine measurement with a discriminating control. This result is L2 on four counts.

\begin{itemize}
\item The g-factor is \emph{read out of the spectrum}, never assigned.
\item The spin gap is \emph{structural}, forced by the sites' Clifford algebra through $H^2$.
\item There is a \emph{clean control}, the spinless scalar, that lacks the gap and the moment.
\item The number is \emph{robust} across fields and masses.
\end{itemize}

The base stays discrete and deterministic throughout. The Landau spectrum is the spectrum of a definite operator on a definite mesh. No randomness, no fitting, no hand-set magnetic response enters at any step.

\subsection{What this does and does not claim}

The claim is exact. The substrate's spinor has the \emph{right magnetic moment}, $g = 2$, derived from its own Dirac structure.

The claim stops there. It says nothing about the value of the coupling that sets how strong the field is. A separate experiment, at depth L1, scans that coupling and finds that the bare knit treats the gauge coupling as a free multiplicative scale. The radiated field follows a clean square law in the coupling with no special or critical value anywhere, vanishing only at zero coupling. A pure power law with no distinguished point means the coupling enters purely as an overall scale, so nothing in the knit fixes its magnitude.

\begin{principle}[Form fixed, strength free]
The knit fixes the \emph{form} of the magnetism, $g = 2$. It does not fix the \emph{strength} of the coupling. These are separate questions. Both are answered honestly, the first with a measurement, the second with a clean negative.
\end{principle}

The two endpoints are different in kind. The form of the moment is forced by the Clifford algebra and comes out as a hard number. The value of the coupling would need an extra principle the bare knit does not contain, a renormalization-group fixed point, anomaly matching, or the full charged-particle spectrum with its running. None of that is present in the knit, and the section claims none of it.

\subsection{Why it matters}

The magnetic moment of the electron is among the most precisely tested numbers in all of physics. A substrate that claims to carry matter cannot leave $g$ to be set by fiat. It must produce $g = 2$ from its own structure, without instruction.

This one does. The chain is short and each link is its own section.

\begin{itemize}
\item The sites carry the spinor (Section~17, Spin and Fermions).
\item The spinor propagates as a Dirac particle (Section~18, Dirac and Lorentz).
\item The spinor has a measured magnetic moment, $g = 2$ (this section).
\end{itemize}

\begin{principle}[The moment is the sites turning]
The magnetic moment is not a setting. It is the sites of a dock turning in a field, and they turn at exactly the rate a fermion should. The number $2$ is the substrate confessing its own Dirac structure.
\end{principle}

The partner measurement, the mass mechanism, is taken up in Section~20, Mass and Chirality, where the $8s$ and $8c$ chiralities couple to give the Dirac mass with the same relativistic form. Together the two close the matter sector on the substrate. The spinor has the right moment and the right mass, both derived rather than assumed, with the values of the coupling and the Yukawa parameter left precisely where the substrate leaves them, open.
